\section{The coefficient-three pure two-power family on original rings}
\label{sec:coefficient-three-high-two}

\paragraph{The coefficient-three case uses a smaller field.}
The theorem AC-G175 is restricted to coefficient primes at least five.
We now treat coefficient three for the pure two-power labels.
The auxiliary field has conductor $3\cdot2^c$, so its degree is prime
to three.  No conductor-seven character or conductor-seven cyclotomic
factor is introduced.  This is a separate class-field and local proof,
not a substitution of three into the field in AC-G175.

Put $K=\mathbb Q(\omega)=\mathbb Q(\mu_3)$, and let $c_3$ denote the
three-adic cyclotomic character, with Teichmuller part $\omega_3$.
Its restriction $\omega_3|_{G_K}$ is identically one.
For a discrete three-primary module $D$, use the full strict condition
\[
 \begin{split}
 \mathcal S_3(K,D)=\ker\bigl(&H^1_{\rm cts}(G_K,D)\longrightarrow
 H^1(G_{K_3},D)\\
 &\times\prod_{v\nmid3,\infty}H^1(I_v,D)\bigr),
 \end{split}\tag{G177.1}
\]
where $K_3=\mathbb Q_3(\mu_3)$ is the unique completion above three.

\begin{theorem}[AC-G177: high two-power strict zero at coefficient three]
\label{AC-G177}
Let $\delta$ be any primitive Dirichlet character of exact conductor
$2^c$, with $c\geq8$, restricted to $G_K$.
Let $E/\mathbb Q_3$ be any finite extension whose original ring of
integers $O$ contains the values of $\delta$.
For every $e\geq1$, $s\in\mathbb Z$, and continuous character
$\kappa:G_K\to1+\mathfrak m_O$ unramified away three, put
\[
 D=(E/O)(\delta\,\omega_3^{-s}\kappa^{-1}c_3^{-e}).
 \tag{G177.2}
\]
Then $\mathcal S_3(K,D)=0$.
The residual strict group is zero over every finite residue field
containing the reduced character values.  There is no weight or
parity exception and no restriction on the absolute ramification
index of the original $O$.

For the actual host, take $s=0$ and $\kappa=1$.
Every such label with $c\leq h+2$ is an actual pure finite label,
with trivial cubic, conductor-seven and auxiliary quadratic factors.
The number of labels supplied at page $h$ is exactly
\[
 C_3(h)=\max\{0,2^{h+1}-64\}.
 \tag{G177.3}
\]
Thus all $C_3(h)(T-1)$ original label--weight pairs with
$1\leq e<T$ have strict group zero at coefficient three.
\end{theorem}

\begin{proof}
\proofstep{1. A fixed-field nonvanishing bound at three.}
Write uniquely
\[
 \delta=\chi_{-4}^{\,a}\psi,\qquad
 a\in\{0,1\},\qquad\psi(-1)=1.
 \tag{G177.4}
\]
Since $c\geq8$, the even character $\psi$ is still primitive of
exact conductor $2^c$.
The fixed factors needed below are
\[
 \theta=\chi_{-4}^{\,a}\chi_{-3}^{\,b},\qquad b\in\{0,1\}.
 \tag{G177.5}
\]
Only the odd factors are used in the special-value argument.
Their primitive conductors divide twelve, and their reductions
belong to $k=\mathbb F_9$, which contains $\mu_4$.
Consequently the notation of (G173.1) gives
\[
 n_0=v_2(9-1)=3,\qquad
 L=\operatorname{lcm}(f_\theta,8)\leq24.
 \tag{G177.6}
\]
For $c\geq8$ all three hypotheses of the general effective
nonvanishing theorem (G173.2) hold:
\[
 c\geq2n_0=6,\qquad v_2(4f_\theta)\leq4<c,\qquad
                  2^{c-n_0}\geq32>23\geq L-1.
 \tag{G177.7}
\]
Therefore every odd product $\theta\psi$ in (G177.5), and its
inverse, has three-adic unit value at zero.
This application includes integrality; no division by a possibly
three-divisible Bernoulli denominator is assumed.
The fixed field $\mathbb F_9$ is independent of the growing values
of $\psi$ and is not the original residue field of $O$.
The complete source of the general theorem remains Sinnott,
\emph{On a theorem of L.~Washington}, Lemma~3.6, printed
pp.~215--217, and Proposition~4.6, printed pp.~220--222, in the
exact effective specialization proved in AC-G173.

\proofstep{2. Odd class sectors in the actual prime-to-three field.}
Set
\[
 H=\mathbb Q(\zeta_{3\cdot2^c}),\qquad
 G=\operatorname{Gal}(H/\mathbb Q),\qquad
 J=\operatorname{Gal}(H/K),\qquad A=\operatorname{Cl}(H)\{3\}.
 \tag{G177.8}
\]
Here $[H:\mathbb Q]=2^c$ and $[H:K]=2^{c-1}$.
After a finite unramified extension of the three-adic coefficient
ring, all characters of $G$ have integral idempotents.
Let $\nu$ be either of the rational extensions of a high two-power
character under consideration, and suppose $\nu$ is odd.
It has the form (G177.5) times an even primitive $\psi$ or
$\psi^{-1}$ and has primitive conductor $f_\nu\mid3\cdot2^c$.
Work first in $H_\nu=\mathbb Q(\zeta_{f_\nu})$.

Choose a positive integer $b_0$ by the Chinese remainder theorem
such that $b_0\equiv5\pmod {2^c}$ and $b_0\equiv1\pmod3$.
The integral Stickelberger theorem, in the convention of Sharifi,
Definition~4.5.1, Lemma~4.5.2 and Theorem~4.5.6, says that
\[
 (b_0-\sigma_{b_0})\Theta_{f_\nu},\qquad
 \Theta_f=\frac1f\sum_{\substack{1\leq u\leq f\\(u,f)=1}}
                                      u\sigma_u^{-1}
 \tag{G177.9}
\]
annihilates the class group of $H_\nu$.
On its $\nu$ component this element acts by
\[
             -(b_0-\nu(b_0))L(0,\nu^{-1}).
 \tag{G177.10}
\]
The second factor is a unit by Step~1.
The first reduces to $1-\overline\psi(5)$, or its inverse-character
version.  Since $5$ generates $1+4\mathbb Z/2^c\mathbb Z$,
$\psi(5)$ has exact order $2^{c-2}$, preserved by reduction at
three.  This factor is therefore a unit too, and the $\nu$
component of the three-class group of $H_\nu$ is zero.

For $V=\operatorname{Gal}(H/H_\nu)$, ideal extension $j$ and norm
$N$ satisfy $Nj=|V|$ and $jN=\sum_{v\in V}v$.
The order of $V$ is a power of two and $\nu$ is trivial on $V$.
The same component is therefore zero in $A$.
This primitive-conductor calculation makes no assumption about
an omitted imprimitive Euler factor.

\proofstep{3. Reflection and its local root-of-unity kernel.}
The field $H$ is CM, contains $\mu_3$, and has degree prime to three.
Let $B\subset H^\times$ be the radicals generating everywhere
unramified Kummer extensions of exponent three.
The class-field/Kummer construction, with $(x)=\mathfrak a^3$,
gives the equivariant sequence
\[
 0\longrightarrow
 (B\cap\mathcal O_H^\times)/\mathcal O_H^{\times3}
 \longrightarrow\operatorname{Hom}(A/3A,\mu_3)
 \longrightarrow A[3].
 \tag{G177.11}
\]
There is no surjectivity assertion at the final term.
We use the exact sequence and equivariance of Sharifi,
Lemma~1.4.6, Proposition~1.3.2 and Theorem~4.4.2, in their
previously audited complete proof scopes.

The odd part of units modulo torsion is zero in a CM field.
The only possible odd three-primary root-of-unity line in
$\mathcal O_H^\times/\mathcal O_H^{\times3}$ is $\mu_3$,
of rational character $\omega_3=\chi_{-3}$.
The conductor $2^c$ extension is unramified at three, so every
completion of $H$ above three has ramification index two over
$\mathbb Q_3$.  Adjoining $\mu_9$ adds relative ramification
index three, unchanged by the unramified factor.
Thus a nonzero $\mu_3$ class cannot belong to the unramified
radical kernel in (G177.11).  Roots of unity of order prime to
three are already cubes in $H$.  The whole odd kernel is zero.

For an even character $\nu$, the component of the middle term
dual to $A(\nu)$ has character $\omega_3\nu^{-1}$.
It injects into $A[3](\omega_3\nu^{-1})$.
That reflected character is odd, has the same exact two-primary
conductor, and is one of the characters killed in Step~2.
Hence $A(\nu)=0$ as well.
Every class component needed for either rational extension of
our original high two-power character is zero, of either parity.
This is not a computation of the entire class group of $H$.

\proofstep{4. Residual strict descent over $K$.}
Let $k'$ be any finite residue field containing the values of
$\overline\delta$.
Because $\omega_3|_{G_K}=1$ and $\kappa$ reduces to one, the
residual character of (G177.2) is simply $\alpha=\overline\delta$,
independently of $e$ and $s$.
Restriction to $G_H$ is injective in degree one since
$[H:K]$ is prime to three, and the coefficient is trivial there.
For a restricted cocycle $z$,
\[
 z(\tau g\tau^{-1})=\alpha(\tau)z(g)
                  \qquad(g\in G_H,\ \tau\in G_K).
 \tag{G177.12}
\]
Strictness kills all inertia and the full decomposition groups
above three.  The resulting homomorphism factors through a
quotient of the ordinary class group, yielding
\[
 \mathcal S_3(K,k'(\alpha))\lhook\joinrel\longrightarrow
          \operatorname{Hom}_{\mathbb F_3[J]}(A/3A,k'(\alpha)).
 \tag{G177.13}
\]
The selected character is $\alpha$, not its inverse.
Its two rational extensions differ by $\chi_{-3}$, and both
class components were killed in Steps~2--3.
Extend $k'$ faithfully to a finite splitting field if needed;
the integral idempotents show the displayed Hom is zero there,
and faithful descent shows it is zero over $k'$.

\proofstep{5. The ramified local base and the original uniformizer.}
The local base is $K_3=\mathbb Q_3(\mu_3)$, ramified quadratic
with residue field $\mathbb F_3$.
The finite character $\delta$ is unramified at three.
An arithmetic Frobenius lift in $G_{K_3}$ therefore has
$\delta$-value $\delta(3)$, with exponent one, not two.
Since $3=-1\cdot(-3)$, $-3\in1+4\mathbb Z_2$ and
$v_2((-3)-1)=2$, the element $-3$ generates
$1+4\mathbb Z/2^c\mathbb Z$.
Thus $\psi(-3)$ has order $2^{c-2}$.
Multiplying by the possible parity sign leaves that exact order
unchanged, since it is greater than two.  Consequently
\[
                 \operatorname{ord}\delta(3)=2^{c-2}>2.
 \tag{G177.14}
\]
Reduction at three preserves this order.
The cyclotomic character modulo three is trivial on all of $G_K$;
so are $\omega_3$ and the reduction of $\kappa$.
The scalar of $D$ at this Frobenius lift minus one is therefore
an $O$-unit.  It follows that $D^{G_{K_3}}=0$, and hence also
$D^{G_K}=0$, for every weight and every allowed $\kappa$.

At each finite place away three, the cyclotomic and $\kappa$
factors are unramified, and the remaining inertia action has
finite two-power order, prime to three.
Its invariants in $D$ are either all of $D$ or zero: every
nontrivial such root of unity differs from one by an $O$-unit.
These invariants are divisible by an original uniformizer
$\varpi$ in either case.
Now $D$ is $\varpi$-divisible and
$D[\varpi]=k'(\overline\delta)$.
The long exact cohomology sequences of
$0\to D[\varpi]\to D\xrightarrow{\varpi}D\to0$, globally,
at $G_{K_3}$ and at every away inertia group, show that every
strict class killed by $\varpi$ lifts to a strict residual class.
Indeed the local maps from residual degree-one cohomology are
injective by the preceding invariant calculation.
Step~4 proves that residual class is zero.
Every continuous cocycle from a profinite group to the discrete
module $D$ has finite image and is killed by a power of $\varpi$.
Thus $\mathcal S_3(K,D)=0$ on the original, arbitrarily ramified ring.

\proofstep{6. Literal host labels and their count.}
The actual host (FG.3) contains $\mathbb Q(\zeta_{2^{h+2}})$.
Every stated $\delta$ with $c\leq h+2$ factors through it, and
distinct such characters remain distinct on $G_K$, since that
two-power field is unramified at three while $K/\mathbb Q$ ramifies.
Their orders divide $2^h$, hence the actual exponent $3\cdot2^h$;
their values lie in the original coefficient ring $O_3$ of (FC.1).
The coefficient in (FC.2) is exactly (G177.2) with $s=0,\kappa=1$.
Inflation from the prescribed ramification set into $G_K$ is
injective and respects the same full local three condition and
all away inertia conditions; outside that set an inflated cocycle
already kills inertia.  Hence its strict group is zero as well.

There are $2^{c-2}$ primitive characters of conductor $2^c$.
Summing from $c=8$ to $c=h+2$ gives
\[
 \sum_{c=8}^{h+2}2^{c-2}=2^{h+1}-2^6
 \quad\text{if }h\geq6,
 \tag{G177.15}
\]
and zero for an empty range.  This proves (G177.3).
The labels are nontrivial, so they are disjoint from the previous
coefficient-three trivial-label row; coefficient-prime separation
also makes them disjoint from every coefficient-at-least-five row.
\end{proof}

\paragraph{Remaining numerical comparison and geometric scope.}
The result is pointwise, independent of the prescribed $q,d$, of
the window and of the weight cutoff.  The passage from strict zero
to a removed positive-index contribution still uses (EU.2) and the
single nonnegative full-union (FL.4) bound on the same window.
No individual index zero is claimed.
The theorem does not cover coefficient two, any conductor-seven
factor, a growing auxiliary quadratic twist or a nontrivial cubic
label, and does not prove the all-prime numerical remainder.
The mixed realization $\mathcal R_{\rm cub}$, H04.4 and cubic RH
remain open.  The finite-etale character and $\mu_3$-torsor language
is the already audited SGA~1, Expos\'e~V, sections4 and7 and
Stacks~03PK dictionary; no new archival attribution is made.
